%----------------------------------------------------------------------------------------
%	STAGE DESCRIPTION
%----------------------------------------------------------------------------------------
\section*{Scopo dello stage}
% Personalizzare inserendo lo scopo dello stage, cioè una breve descrizione
Lo scopo di questo progetto di stage è progettare e sviluppare una piattaforma web collaborativa dedicata alla condivisione e revisione di codice, con l’obiettivo di facilitare l’interazione tra sviluppatori di diversi livelli di esperienza.
La piattaforma intende fornire un ambiente strutturato in cui gli utenti possano pubblicare snippet di codice da loro realizzati, ricevere feedback qualificati e migliorare le proprie competenze attraverso un confronto diretto e continuo. Il progetto mira inoltre a esplorare soluzioni tecnologiche moderne, integrando un'architettura frontend sviluppata in Angular e un backend realizzato in Spring.

\vspace{1em}

Lo studente avrà il compito di analizzare i requisiti funzionali e non funzionali della piattaforma, progettare l'architettura del sistema e implementarne le principali funzionalità. Sarà inoltre responsabile della scelta e integrazione degli strumenti di sviluppo e della validazione del sistema tramite test e valutazioni delle prestazioni. Durante lo svolgimento del progetto, lo studente dovrà adottare buone pratiche di sviluppo software, documentare il lavoro svolto, comunicare in maniera attiva con l'azienda e il tutor aziendale e dimostrare la capacità di affrontare problematiche reali in un contesto applicativo concreto, comprendendo in maniera profonda le motivazioni dietro alle scelte architetturali compiute.

